\documentclass{name}
\usepackage{lipsum}
\usepackage{multicol}
\usepackage{listings}
\usepackage{longtable}
\usepackage{amsmath}
\usepackage{mathtools}
\usepackage{amssymb}
\usepackage{booktabs}
\usepackage{array}
\usepackage{caption}
\usepackage{subcaption}
\usepackage{float}
\usepackage{setspace}
\usepackage{stmaryrd}
\usepackage[backend=bibtex,sorting=none]{biblatex}
\addbibresource{bibliographie.bib}
\usepackage{csquotes}
\usepackage[english]{babel}
\usepackage{fancyhdr}
\usepackage{lmodern}
\usepackage{xcolor}
\usepackage{algorithm}
\usepackage{algorithmic}
\usepackage{awesomebox}
\usepackage{todonotes}
\usepackage{upquote}

\newtheorem{theorem}{Theorem}
\newtheorem{remark}{Remark}
\newtheorem{definition}{Definition}

\newcounter{pagezero}

\begin{document}

%----------- Informations du rapport ---------
\UE{S8 Machine Learning}
\sujet{Kaggle Project: Textile Classification \& AI Safety}
\titre{TILDA Textile Classifier and AI Safety}

\eleves{Marin \textsc{Decanini}\\Nolan \textsc{Carlisi}}

\enseignant{Pierre \textsc{Weiss}\\Université de Toulouse\\ Senior CNRS researcher at IRIT}

%----------- Initialisation -------------------
\fairemarges %Afficher les marges
\fairepagedegarde %Créer la page de garde

\newpage
\setcounter{pagezero}{0}

\begin{center}
	\begin{minipage}{13cm}
	    \begin{center}
	        \section*{Abstract}
	    \end{center}
	    This report presents our implementation and experimental analysis for the S8 Machine Learning Kaggle project on the TILDA textile dataset. 
	    We focus on two core components: first, training state-of-the-art Convolutional Neural Networks (CNNs) for classification; and second, analyzing machine learning security and AI safety under biased training distributions. 
	    Specifically, we implement LeNet and ResNet architectures, evaluate their performance using mini-batch SGD, and perform systematic hyperparameter tuning. 
	    Furthermore, we investigate model vulnerability to spurious correlations by constructing a biased dataset with a noise-like channel correlated with labels, measuring both classification accuracy and fairness using the Disparate Impact (DI) metric.
	\end{minipage}
\end{center}

\par 
\setcounter{page}{\value{pagezero}}

\newpage
\setcounter{page}{\value{pagezero}}
\tableofcontents

\newpage

\section{Introduction}
% Introduce the report, goals, and structure.

\section{Training State-of-the-Art CNN Models}

\subsection{Image Pre-processing and Data Augmentation}
% Question 2.1: Explain the pre-processing steps and data augmentation strategies.

\subsection{Model Architecture and Selection}
% Question 2.2: Rationale for choosing the networks and precise definition of the models (layers, dimensions, functions).

\subsection{Training and Optimization Methodology}
% Question 2.3: Details on mini-batch SGD, optimization parameters, hardware, and training times.

\subsection{Hyperparameter Tuning and Validation}
% Question 2.4: Validation strategy, hyperparameter optimization process, and performance summary.

\section{Machine Learning and AI Safety}

\subsection{Real-Life Bias Context}
% Question 3.1: Discussion on real-world AI bias examples and their societal/ethical consequences.

\subsection{Biased Dataset Construction}
% Question 3.2: Explanation of the mathematical setup (Bernoulli trial, noise channel epsilon) and Pytorch dataset implementation.

\subsection{Model Bias Evaluation}

\subsubsection{Disparate Impact Analysis}
% Question 3.3: Comparative evaluation of DI metrics between model 1 (unbiased) and model 2 (biased).

\subsubsection{Test Accuracy Comparison}
% Question 3.4: Test accuracy performance of model 1 vs model 2 on the binary classification task.

\subsubsection{Summary of Results and Takeaways}
% Question 3.5: Table summarizing Train/Val/Test accuracy and DI for both models, with conclusions.

\subsection{Bias Study in the Literature (Bonus)}
% Question 3.6: Short essay on "Est-ce qu'on peut faire confiance à mon modèle?" in light of model bias guidelines (EU requirements, CNIL).

\printbibliography

\end{document}
